
tères pour que @@M0@@ soit injectif (resp. surjectif).

\textbf{7.18.2.} Supposons que l'on soit dans l'un des deux cas suivants :
\begin{enumerate}[label=\alph*)]
\item Le corps @@M1@@ est séparablement clos, et @@M2@@ est une extension séparable.
\item Le corps @@M3@@ est infini, et @@M4@@ en est une extension pure.
\end{enumerate}
Alors l'application @@M5@@ est injective.

En effet, utilisant le fait que @@M6@@ est limite inductive filtrante de sous-algèbres de type fini sur @@M7@@, avec des morphismes de transition plats, et la commutation de @@M8@@ aux limites inductives filtrantes, on est ramené dans le cas a) à prouver que si @@M9@@ est une algèbre intègre de type fini sur @@M10@@, dont le corps des fractions @@M11@@ est séparable sur @@M12@@, alors @@M13@@ est injectif. Or les hypothèses faites sur @@M14@@ et sur @@M15@@ impliquent qu'il existe un @@M16@@-homomorphisme @@M17@@; alors le composé @@M18@@ est l'identité, d'où la conclusion. Dans le cas b), on se ramène de même, par passage à la limite, à prouver que si @@M19@@ est une algèbre intègre de type fini sur @@M20@@ telle que son corps des fractions soit une extension pure de @@M21@@, alors @@M22@@ est injectif. Comme @@M23@@ est infini, l'hypothèse implique encore qu'il existe un @@M24@@-homomorphisme @@M25@@, et on achève comme ci-dessus.

\textbf{7.18.3.} Supposons maintenant que @@M26@@ soit un foncteur à valeurs dans la catégorie @@M27@@ des groupes abéliens, et supposons que pour toute extension finie @@M28@@ de @@M29@@, on ait défini un homomorphisme norme
@@M30@@
tel que le composé
@@M31@@
soit la multiplication par @@M32@@. Sous ces conditions, on a :
\begin{enumerate}[label=(\roman*)]
